% =====================================================================
% Fluid Dynamics and IBM Theory: motivation for the immersed boundary method
% =====================================================================
\documentclass[aspectratio=169,11pt]{beamer}

% Shared style, ported from the UF Animal Sciences 2026 talk (ib-tutorials.sty at the repo root).
\makeatletter\def\input@path{{./}{../}}\makeatother
\usepackage{ib-tutorials}

\usepackage[utf8]{inputenc}
\usepackage{graphicx}
\usepackage{amsmath,amssymb,bm}
\usepackage{booktabs,tabularx,array}
\usepackage{xcolor}
\usepackage{tcolorbox}
\tcbuselibrary{skins}
\usepackage{hyperref}

\graphicspath{{../pictures/}{pictures/}}

\newcommand{\xvec}{\bm{x}}
\newcommand{\uvec}{\bm{u}}
\newcommand{\Xvec}{\bm{X}}
\newcommand{\Fvec}{\bm{F}}
\newcommand{\fvec}{\bm{f}}
\newcommand{\Uvec}{\bm{U}}
\newcommand{\jump}[1]{\left[\!\left[#1\right]\!\right]}

% Title-slide image followed by its source credit (see the title frame).
\let\tutorialgraphic\includegraphics
\newcommand{\creditedtitlegraphic}[2][]{%
  \tutorialgraphic[#1]{#2}\par\vspace{2pt}
  {\scriptsize Schematic after the Wikimedia Commons figure \emph{ImmersedBoundary grid} (Immersed boundary method article).\par}}

\title[Immersed boundary method]{Motivation for the immersed boundary method}
\subtitle{Classical Peskin formulation, tradeoffs, and method selection}
\author[Miller]{Laura A. Miller}
\institute[UArizona]{Department of Mathematics, University of Arizona}
\date{Fluid Dynamics and IBM Theory \textbullet\ Module 1}

\begin{document}

% ----- 1. Title -----
% The image credit is kept under the title-slide image: \includegraphics is
% wrapped for this frame only so \maketutorialtitle prints the credit too.
\begin{frame}[plain]
\let\includegraphics\creditedtitlegraphic
\maketutorialtitle[ImmersedBoundary_grid2.png]
\end{frame}

% ----- 2. Definition -----
\begin{frame}{What is the immersed boundary method?}
\begin{columns}[T,onlytextwidth]
\begin{column}{0.53\textwidth}
The immersed boundary (IB) method is a computational framework for fluid--structure interaction in which
\begin{itemize}\itemsep6pt
  \item fluid momentum is described on a fixed \hi{Eulerian} grid;
  \item structure position and mechanics are described with moving \hi{Lagrangian} coordinates;
  \item integral transforms transfer force to the fluid and velocity back to the structure.
\end{itemize}
\end{column}
\begin{column}{0.43\textwidth}
\begin{punchbox}
\centering
\small
\[
\begin{gathered}
\text{structure force}\xrightarrow{\;S\;}\text{fluid force},\\[3pt]
\text{fluid velocity}\xrightarrow{\;J\;}\text{structure motion}.
\end{gathered}
\]
{\scriptsize The spreading operator $S$ and the interpolation operator $J$ are defined below: continuum forms on the ``Continuum force and velocity coupling'' slide, discrete forms $S_h$, $J_h$ on the slide after it.}
\end{punchbox}
\medskip
\small Introduced by Charles S. Peskin for cardiac fluid dynamics, the framework now supports many thin-boundary and volumetric FSI models.
\end{column}
\end{columns}
\end{frame}

% ----- 2. Motivation -----
\begin{frame}{Why use an immersed formulation?}
\begin{columns}[T,onlytextwidth]
\begin{column}{0.48\textwidth}
\textbf{Body-fitted fluid meshes}
\begin{itemize}\itemsep5pt
  \item conform to the moving fluid--structure interface;
  \item can resolve interfacial flow and stress sharply;
  \item may require mesh smoothing or remeshing when motion, deformation, or contact degrades mesh quality.
\end{itemize}
\end{column}
\begin{column}{0.48\textwidth}
\textbf{The immersed alternative}
\begin{itemize}\itemsep5pt
  \item keeps the Eulerian fluid grid fixed;
  \item moves a nonconforming Lagrangian representation of the structure through that grid;
  \item couples the two descriptions by transferring force and velocity.
\end{itemize}
\end{column}
\end{columns}
\medskip
\begin{punchbox}
\small Geometric flexibility comes from the nonconforming fluid and structure descriptions. The Lagrangian structure is still tracked explicitly.
\end{punchbox}
\end{frame}

% ----- 3. Scope and history -----
\begin{frame}{The classical Peskin IB method}
\begin{columns}[T,onlytextwidth]
\begin{column}{0.41\textwidth}
\centering
\includegraphics[width=0.72\linewidth]{ImmersedBoundary_grid.png}

{\scriptsize Blue: fixed Eulerian grid. Orange: moving Lagrangian points.
Schematic after the Wikimedia Commons figure \emph{ImmersedBoundary grid} (Immersed boundary method article).}
\end{column}
\begin{column}{0.55\textwidth}
\begin{itemize}\itemsep3pt
  \item Charles S. Peskin introduced the method to study flow around heart valves in the 1970s.
  \item The classical formulation couples an elastic boundary to a viscous incompressible fluid.
  \item Velocity, pressure, and momentum use an \hi{Eulerian} description. Structure position and elasticity use a \hi{Lagrangian} description.
  \item Later IB extensions treat volumetric solids, active materials, prescribed motion, and constrained rigid bodies.
\end{itemize}
\smallskip
\soft{\small Applications include cardiac flow, cilia and flagella, animal locomotion, flexible beams, and membranes.}
\end{column}
\end{columns}
\vfill
{\scriptsize Peskin, \emph{J. Comput. Phys.} 10 (1972) and 25 (1977); Peskin, \emph{Acta Numerica} 11 (2002); Griffith and Patankar, \emph{Annu. Rev. Fluid Mech.} 52 (2020).}
\end{frame}

% ----- Applications -----
\begin{frame}{Representative applications}
\begin{columns}[T,onlytextwidth]
\begin{column}{0.49\textwidth}
\hi{Biological fluid mechanics}
\begin{itemize}\itemsep5pt
  \item heart valves, chambers, and vascular flow;
  \item cilia- and flagella-driven transport;
  \item swimming by fish, jellyfish, and microorganisms;
  \item flexible airways and tissue--fluid interaction.
\end{itemize}
\end{column}
\begin{column}{0.49\textwidth}
\hi{Engineering and soft matter}
\begin{itemize}\itemsep5pt
  \item flapping flags, membranes, and beams;
  \item flow-induced vibration and mixing;
  \item bio-inspired propulsion and soft robotics;
  \item compliant or porous surfaces, with extensions tailored to the required physics.
\end{itemize}
\end{column}
\end{columns}
\begin{punchbox}
\small The common feature is not a single Reynolds-number regime; it is a moving or deforming interface that benefits from a nonconforming fluid grid.
\end{punchbox}
\end{frame}

% ----- Scope -----
\begin{frame}{What the classical IB method is---and is not}
\begin{columns}[T,onlytextwidth]
\begin{column}{0.49\textwidth}
\textbf{It is}
\begin{itemize}\itemsep6pt
  \item a force-based FSI formulation;
  \item an Eulerian--Lagrangian coupling framework;
  \item a method whose standard discretization uses a regularized delta kernel;
  \item compatible with large structural motion.
\end{itemize}
\end{column}
\begin{column}{0.49\textwidth}
\textbf{It is not}
\begin{itemize}\itemsep6pt
  \item a body-fitted or intrinsically sharp-interface method;
  \item a method that discards elasticity or stress;
  \item automatically topology changing---the structural representation must model any change;
  \item restricted to one constitutive law.
\end{itemize}
\end{column}
\end{columns}
\begin{punchbox}
\small At the continuum level no slip is pointwise. ``Weak enforcement'' describes the regularized discrete interaction, not the singular continuum equations.
\end{punchbox}
\end{frame}

% ----- 4. Continuum coupling -----
\begin{frame}{Continuum force and velocity coupling}
Let $q\in\mathcal U$ label a material point and let $\Xvec(q,t)$ be its position in the fluid domain $\Omega$.
\begin{columns}[T,onlytextwidth]
\begin{column}{0.43\textwidth}
\textbf{Spread structural force to the fluid}
\[
  \fvec(\xvec,t)
  =\int_{\mathcal U}\Fvec(q,t)\,
  \delta\!\left(\xvec-\Xvec(q,t)\right)\,dq.
\]
$\Fvec$ is a force density per unit material measure. The Dirac distribution places that force at the current structure location.
\end{column}
\begin{column}{0.53\textwidth}
\textbf{Interpolate fluid velocity to the structure}
\[
  \Uvec(q,t)=\frac{\partial\Xvec}{\partial t}
  =\int_{\Omega}\uvec(\xvec,t)\,
  \delta\!\left(\xvec-\Xvec(q,t)\right)\,d\xvec.
\]
For continuous $\uvec$, the Dirac identity gives
\[
  \frac{\partial\Xvec}{\partial t}(q,t)
  =\uvec\!\left(\Xvec(q,t),t\right).
\]
\end{column}
\end{columns}
\smallskip
\begin{punchbox}
\small At the continuum level, the kinematic equation imposes pointwise no slip on an impermeable material interface. It is not merely an averaged boundary condition.
\end{punchbox}
{\scriptsize In operator form, $\fvec=S[\Xvec]\,\Fvec$ and $\partial_t\Xvec=J[\Xvec]\,\uvec$ with $J=S^{*}$, the $L^2$ adjoint of $S$. Adapted from Griffith and Patankar (2020), \S2.1.1, Eqs.~1--4, and Peskin (2002).\par}
\end{frame}

% ----- 5. Discrete coupling -----
\begin{frame}{The standard discrete IB coupling}
The numerical method replaces $\delta$ by a regularized kernel $\delta_h$ whose support is a few Eulerian grid cells wide:
\[
  \fvec_h=S_h\Fvec_h,
  \qquad
  \Uvec_h=J_h\uvec_h.
\]
\begin{columns}[T,onlytextwidth]
\begin{column}{0.50\textwidth}
\begin{itemize}\itemsep3pt
  \item $S_h$ spreads Lagrangian forces to nearby grid points.
  \item $J_h$ interpolates grid velocity to Lagrangian interaction points.
  \item Kernel support and moment conditions control smoothness, localization, and grid dependence.
\end{itemize}
\end{column}
\begin{column}{0.46\textwidth}
Choose $J_h$ as the weighted discrete adjoint of $S_h$ so that the transferred power is consistent:
\[
  \left\langle \uvec_h,S_h\Fvec_h\right\rangle_E
  =
  \left\langle J_h\uvec_h,\Fvec_h\right\rangle_L.
\]
\end{column}
\end{columns}
\begin{punchbox}
\small Regularization spreads a singular force layer and associated jump fields over $O(h)$. The Lagrangian boundary itself remains explicitly represented.
\end{punchbox}
\end{frame}

% ----- Adjointness (added worked step) -----
\begin{frame}{Why adjoint operators transfer power consistently}
With Eulerian grid spacing $h$ and Lagrangian quadrature weights $\Delta q_l$, define
\[
  \langle\uvec_h,\fvec_h\rangle_E=\sum_{\xvec}\uvec_h(\xvec)\cdot\fvec_h(\xvec)\,h^3,
  \qquad
  \langle\Uvec_h,\Fvec_h\rangle_L=\sum_{l}\Uvec_l\cdot\Fvec_l\,\Delta q_l.
\]
\begin{itemize}\itemsep4pt
  \item Power the spread force does on the fluid: $P_E=\langle\uvec_h,S_h\Fvec_h\rangle_E$.
  \item Power the structural force does at the interpolated structure velocity: $P_L=\langle J_h\uvec_h,\Fvec_h\rangle_L$.
  \item If $J_h=S_h^{*}$, then $P_E=P_L$ for every $\uvec_h$ and $\Fvec_h$: the transfer step neither creates nor destroys energy.
\end{itemize}
{\scriptsize Here $\langle\cdot,\cdot\rangle_E$ and $\langle\cdot,\cdot\rangle_L$ are the discrete Eulerian and Lagrangian inner products induced by the quadrature rules. See Griffith and Patankar (2020), \S3.3, who write these operators as $S_{\mathrm{IB}}$ and $J_{\mathrm{IB}}$ with $J_{\mathrm{IB}}=S_{\mathrm{IB}}^{*}=h^3M^{-1}S_{\mathrm{IB}}^{\mathsf T}$, where $M$ is the diagonal matrix of Lagrangian quadrature weights (for uniform weights $\Delta q$, $J_{\mathrm{IB}}=(h^3/\Delta q)\,S_{\mathrm{IB}}^{\mathsf T}$).\par}
\end{frame}

% ----- 6. Strengths and limitations -----
\begin{frame}{What regularized IB coupling buys and costs}
\begin{columns}[T,onlytextwidth]
\begin{column}{0.48\textwidth}
\textbf{Strengths}
\begin{itemize}\itemsep3pt
  \item No body-conforming fluid mesh or fluid-grid remeshing
  \item Large structural motion and deformation
  \item Thin immersed boundaries and transient contact
  \item Efficient Cartesian-grid solvers and adaptive mesh refinement
  \item Fiber, surface, and volumetric structural models
\end{itemize}
\end{column}
\begin{column}{0.48\textwidth}
\textbf{Numerical challenges}
\begin{itemize}\itemsep3pt
  \item Pressure and traction jumps are smoothed at finite $h$
  \item Interfacial forces and stresses may converge slowly
  \item Coarse interaction-point spacing can cause leakage or volume error
  \item Stiff elastic forces can severely restrict explicit time steps
\end{itemize}
\end{column}
\end{columns}
\medskip
\soft{\small AMR targets spatial resolution. Implicit or semi-implicit coupling targets stiffness. Sharp-interface corrections target jump accuracy. Finite-element structure models broaden constitutive and geometric choices.}
\end{frame}

% ----- Accuracy -----
\begin{frame}{Accuracy and regularized boundary coupling}
\begin{columns}[T,onlytextwidth]
\begin{column}{0.50\textwidth}
At finite grid spacing, \(\delta_h\) averages information across a stencil several cells wide.
\begin{itemize}\itemsep5pt
  \item The interpolated velocity need not equal the pointwise trace of a sharply resolved field.
  \item Pressure and stress jumps are represented over an \(O(h)\) layer.
  \item Kernel moments, support, and smoothness influence conservation and grid dependence.
\end{itemize}
\end{column}
\begin{column}{0.48\textwidth}
\begin{punchbox}
\small Accuracy depends jointly on
\[
h,\qquad \Delta s\ \text{or quadrature density},\qquad \Delta t,
\]
the fluid solver, the structural discretization, and solution regularity.
\end{punchbox}
\medskip
\small A refinement study must vary every resolution that controls the coupled approximation.
\end{column}
\end{columns}
\end{frame}

% ----- Failure modes -----
\begin{frame}{Common failure modes and practical checks}
\begin{columns}[T,onlytextwidth]
\begin{column}{0.49\textwidth}
\textbf{Interaction resolution}
\begin{itemize}\itemsep5pt
  \item Interaction points too far apart can leak.
  \item A very fine structural mesh may add stiffness or poor conditioning without improving the fluid resolution.
  \item Different spreading and interpolation kernels break the usual adjoint power identity.
\end{itemize}
\end{column}
\begin{column}{0.49\textwidth}
\textbf{Time integration and verification}
\begin{itemize}\itemsep5pt
  \item Explicit elastic forces can impose severe time-step limits.
  \item Monitor volume, deformation, net force, and transferred power separately.
  \item Refine Eulerian, Lagrangian, quadrature, and temporal resolutions systematically.
\end{itemize}
\end{column}
\end{columns}
\begin{punchbox}
\small Many apparent ``IB failures'' are mismatches among the grid, interaction points, kernel, and time integrator rather than failures of the continuum formulation.
\end{punchbox}
{\scriptsize With nodal quadrature, a structure leaks when its Lagrangian nodes become too far apart relative to $h$ (Peskin 2002). Non-nodal (e.g.\ Gaussian) quadrature lets the structural mesh be coarse provided the net of quadrature points is dense enough (Griffith and Luo 2017; Griffith and Patankar 2020, \S3.3).\par}
\end{frame}

% ----- 7. Method selection -----
\begin{frame}{When is classical IB a good choice?}
\begin{columns}[T,onlytextwidth]
\begin{column}{0.48\textwidth}
\textbf{Strong fit}
\begin{itemize}\itemsep3pt
  \item Large or complicated structural motion
  \item Thin, flexible, or actively deforming structures
  \item Transient contact without a body-fitted fluid mesh
  \item Quantities of interest that tolerate regularized interfacial fields
\end{itemize}
\end{column}
\begin{column}{0.48\textwidth}
\textbf{Requires special care}
\begin{itemize}\itemsep3pt
  \item Pointwise tractions or high-order interface stresses
  \item Thin high-Reynolds-number boundary layers
  \item Nearly rigid constraints or severe stiffness
  \item Tearing, merging, or reconnection, which require an explicit topology model
\end{itemize}
\end{column}
\end{columns}
\begin{punchbox}
\small High Reynolds number is not a formal exclusion. Published immersed FSI examples reach $\mathrm{Re}\approx 2\times10^4$, but thin shear layers and turbulence still require adequate resolution and modeling.
\end{punchbox}
{\tiny Griffith and Patankar, \emph{Annu. Rev. Fluid Mech.} 52 (2020); the $\mathrm{Re}\approx2\times10^4$ example is the bioprosthetic heart-valve application of their \S5.3.1.}
\end{frame}

% ----- 8. Numerical practice -----
\begin{frame}{Discretization and convergence checklist}
\begin{columns}[T,onlytextwidth]
\begin{column}{0.48\textwidth}
\textbf{Interaction geometry}
\begin{itemize}\itemsep3pt
  \item Resolve the structure with enough Lagrangian interaction points or quadrature points relative to $h$.
  \item Coarse nodal coupling can leak. Excessively fine structural meshes add cost and may worsen stiffness or conditioning.
  \item Preserve weighted adjointness between spreading and interpolation.
\end{itemize}
\end{column}
\begin{column}{0.48\textwidth}
\textbf{Verification}
\begin{itemize}\itemsep3pt
  \item Refine $h$, the Lagrangian discretization or quadrature, and $\Delta t$ consistently.
  \item Check velocity, deformation, volume, loads, and energy transfer separately.
  \item Use an implicit or semi-implicit treatment when elastic stiffness controls stability.
\end{itemize}
\end{column}
\end{columns}
\begin{punchbox}
\small If the support of $\delta_h$ scales with $h$ and the coupled scheme is stable and consistent, refinement recovers the singular IB formulation. Convergence rates depend on solution regularity; standard regularized methods are often only first order near stress jumps.
\end{punchbox}
{\scriptsize Peskin (2002); Griffith and Patankar (2020); Griffith and Luo (2017).}
\end{frame}

% ----- Body-fitted methods -----
\begin{frame}{Body-fitted FSI: why conforming meshes are attractive}
\begin{columns}[T,onlytextwidth]
\begin{column}{0.50\textwidth}
A body-fitted fluid mesh conforms to the moving interface.
\begin{itemize}\itemsep6pt
  \item Boundary conditions and traction continuity can be imposed on the represented surface.
  \item Boundary layers and interfacial stresses can be resolved without regularized smearing.
  \item ALE, moving-mesh finite element, and moving-mesh finite volume schemes are common examples.
\end{itemize}
\end{column}
\begin{column}{0.46\textwidth}
\begin{punchbox}
\centering
\textbf{Geometric bargain}\\[5pt]
sharp, conforming interface
\[
\Updownarrow
\]
fluid-mesh motion and quality control
\end{punchbox}
\end{column}
\end{columns}
\end{frame}

% ----- Body-fitted costs -----
\begin{frame}{Body-fitted FSI: cost under large motion}
\begin{columns}[T,onlytextwidth]
\begin{column}{0.50\textwidth}
\textbf{Typical numerical work}
\begin{itemize}\itemsep6pt
  \item solve for mesh motion or smoothing;
  \item detect distortion and remesh when needed;
  \item transfer fields between successive meshes;
  \item maintain coupled fluid--solid interface data structures.
\end{itemize}
\end{column}
\begin{column}{0.48\textwidth}
\textbf{Best fit}
\begin{itemize}\itemsep6pt
  \item interface accuracy is central;
  \item motion remains compatible with mesh quality;
  \item the added geometric machinery is justified by the quantities of interest.
\end{itemize}
\end{column}
\end{columns}
\begin{punchbox}
\small Large deformation is possible in sophisticated body-fitted schemes; it is the mesh-management burden, not a universal deformation cutoff, that distinguishes the tradeoff.
\end{punchbox}
\end{frame}

% ----- 9. Comparison -----
\begin{frame}{Comparison with other FSI approaches}
\footnotesize
\renewcommand{\arraystretch}{1.28}
\begin{tabularx}{\textwidth}{@{}>{\raggedright\arraybackslash}p{0.20\textwidth}>{\raggedright\arraybackslash}p{0.34\textwidth}>{\raggedright\arraybackslash}X@{}}
\toprule
\textbf{Method family} & \textbf{Fluid and interface treatment} & \textbf{Typical numerical challenge} \\
\midrule
Regularized IB
& Fixed, nonconforming grid; $\delta_h$ spreading and interpolation
& Smeared jump fields, interfacial accuracy, leakage, and stiffness \\
\addlinespace
ALE / body-fitted
& Moving, boundary-conforming fluid mesh
& Mesh quality, mesh motion, and possible remeshing under large motion \\
\addlinespace
Immersed interface / sharp immersed
& Fixed grid; interface-aware jump corrections
& Geometric reconstruction and jump-dependent stencils or solvers \\
\addlinespace
Fictitious domain / penalization
& Fixed background grid; distributed constraints or penalties
& Conditioning, penalty or modeling error, and method-specific coupling \\
\bottomrule
\end{tabularx}
\medskip
\begin{punchbox}
\small The deformation a method can handle depends on its interface representation, mesh strategy, solver, and implementation. Small, moderate, and large are not reliable method-family labels.
\end{punchbox}
\end{frame}

% ----- 10. IIM -----
\begin{frame}{The immersed interface method}
\small
\begin{columns}[T,onlytextwidth]
\begin{column}{0.52\textwidth}
The immersed interface method (IIM) keeps a fixed Eulerian grid and incorporates analytically derived jump conditions into nearby stencils.
\begin{itemize}\itemsep4pt
      \item It avoids regularized smearing of jump fields.
  \item It can improve near-interface velocity, pressure, deformation, and traction accuracy.
  \item It adds geometry reconstruction and jump-dependent stencils or solvers.
\end{itemize}
\end{column}
\begin{column}{0.45\textwidth}
Define $\jump{a}=a^+-a^-$ and take $\bm n$ from $\Omega^-$ to $\Omega^+$. If $\bm f_\Gamma$ is the interface-on-fluid force in $+\bm f_\Gamma\delta_\Gamma$, then
\[
  \jump{\uvec}=0,
  \qquad
  \jump{\bm\sigma}\bm n=-\bm f_\Gamma.
\]
Traction balance determines jumps in pressure and velocity derivatives. Pressure may be continuous or discontinuous, depending on the interfacial force.
\end{column}
\end{columns}
\begin{punchbox}
\small Sharp means resolving jumps without regularized smearing; it does not mean the computed solution is exact. Large or topology-changing motion still depends on the interface and geometry algorithms.
\end{punchbox}
{\scriptsize Li and Lai, \emph{J. Comput. Phys.} 171 (2001); Griffith and Patankar (2020).}
\end{frame}

% ----- Jump-condition sign check (added worked step) -----
\begin{frame}{Worked step: the sign of the stress jump}
Consider a steady one-dimensional balance with a point force at an interface $x=0$, normal $\bm n=+\bm e_x$ pointing from $\Omega^-=\{x<0\}$ to $\Omega^+=\{x>0\}$:
\[
  0=\frac{\partial\sigma}{\partial x}+f_\Gamma\,\delta(x).
\]
Integrate across $(-\epsilon,\epsilon)$ and let $\epsilon\to0$:
\[
  \sigma(0^+)-\sigma(0^-)+f_\Gamma=0
  \qquad\Longrightarrow\qquad
  \boxed{\jump{\sigma}=-f_\Gamma}.
\]
In several dimensions the same pillbox argument gives $\jump{\bm\sigma}\bm n=-\bm f_\Gamma$. With $\bm\sigma=-p\bm I+\mu(\nabla\uvec+\nabla\uvec^{\mathsf T})$, the normal component yields the pressure jump and the tangential components yield jumps in velocity derivatives.
\begin{punchbox}
\small If the force pushes the fluid toward $\Omega^+$ ($f_\Gamma>0$), the stress on the $+$ side is lower; for a fluid at rest ($\sigma=-p$) this means $p(0^+)=p(0^-)+f_\Gamma$, so the force compresses the fluid ahead of it.
\end{punchbox}
{\scriptsize Pillbox derivation in the spirit of Peskin (2002) and Li and Lai (2001).\par}
\end{frame}

% ----- IIM tradeoffs -----
\begin{frame}{Immersed interface methods: strengths and tradeoffs}
\begin{columns}[T,onlytextwidth]
\begin{column}{0.49\textwidth}
\textbf{Strengths}
\begin{itemize}\itemsep6pt
  \item fixed background grid;
  \item sharp representation of pressure and stress jumps;
  \item potentially higher near-interface accuracy;
  \item no regularized force layer in the target continuum description.
\end{itemize}
\end{column}
\begin{column}{0.49\textwidth}
\textbf{Added requirements}
\begin{itemize}\itemsep6pt
  \item interface geometry and normals;
  \item jump conditions consistent with the physical model;
  \item corrected stencils or specialized solvers;
  \item robust handling of large motion and topology changes.
\end{itemize}
\end{column}
\end{columns}
\begin{punchbox}
\small IIM prioritizes sharp jump accuracy; regularized IB prioritizes a simple, robust transfer mechanism. The right choice follows the observable and geometry.
\end{punchbox}
\end{frame}

% ----- 11. Software and tutorials -----
\begin{frame}{Open-source software and tutorials}
\begin{columns}[T,onlytextwidth]
\begin{column}{0.48\textwidth}
\textbf{IBAMR}
\begin{itemize}\itemsep4pt
  \item Distributed-memory C++ framework
  \item Cartesian adaptive mesh refinement
  \item Examples for classical IB, IBFE, constrained IB, and IIM in two and three dimensions
\end{itemize}
\href{https://ibamr.github.io/}{\texttt{ibamr.github.io}}

{\scriptsize Griffith et al., \emph{J. Comput. Phys.} 223 (2007).\par}
\end{column}
\begin{column}{0.48\textwidth}
\textbf{IB2d}
\begin{itemize}\itemsep4pt
  \item MATLAB and Python implementations in two dimensions
  \item More than 75 examples for teaching and research
  \item Introductory video tutorials linked from the repository
\end{itemize}
\href{https://github.com/nickabattista/IB2d}{\texttt{github.com/nickabattista/IB2d}}

{\scriptsize Battista, Strickland, and Miller, \emph{Bioinspir. Biomim.} 12 (2017).\par}
\end{column}
\end{columns}
\medskip
\begin{punchbox}
\small \href{https://www.youtube.com/watch?v=PJyQA0vwbgU}{Nick Battista's introduction to the IB method} is the first tutorial linked by IB2d. Course-specific \href{https://drive.google.com/drive/folders/1G2bkEsv2Jz_R7165a5BqPbW4dU2hXBHb?usp=sharing}{lectures by Matea Santiago} may require Google Drive access.
\end{punchbox}
\end{frame}

% ----- Open-source value -----
\begin{frame}{Why the open-source ecosystem matters}
\begin{columns}[T,onlytextwidth]
\begin{column}{0.49\textwidth}
\textbf{For learning and verification}
\begin{itemize}\itemsep6pt
  \item inspect the actual spreading and interpolation operators;
  \item reproduce benchmark problems;
  \item compare kernels, solvers, and time integrators;
  \item trace continuum notation into implementation.
\end{itemize}
\end{column}
\begin{column}{0.49\textwidth}
\textbf{For research}
\begin{itemize}\itemsep6pt
  \item prototype new constitutive or active-force models;
  \item use tested parallel solvers and AMR infrastructure;
  \item document algorithmic choices for reproducibility;
  \item separate model development from low-level solver development.
\end{itemize}
\end{column}
\end{columns}
\begin{punchbox}
\small Software availability is a practical advantage, not a mathematical accuracy claim. Validation is still problem- and implementation-specific.
\end{punchbox}
\end{frame}

% ----- 12. References -----
\begin{frame}{Primary references}
\small
\begin{itemize}\itemsep7pt
  \item C. S. Peskin, \emph{The immersed boundary method}, \emph{Acta Numerica} 11 (2002), 479--517. \href{https://doi.org/10.1017/S0962492902000077}{doi:10.1017/S0962492902000077}.
  \item B. E. Griffith and N. A. Patankar, \emph{Immersed methods for fluid--structure interaction}, \emph{Annual Review of Fluid Mechanics} 52 (2020), 421--448. \href{https://doi.org/10.1146/annurev-fluid-010719-060228}{doi:10.1146/annurev-fluid-010719-060228}.
  \item B. E. Griffith and X. Y. Luo, \emph{Hybrid finite difference/finite element immersed boundary method}, \emph{International Journal for Numerical Methods in Biomedical Engineering} 33 (2017), e2888. \href{https://doi.org/10.1002/cnm.2888}{doi:10.1002/cnm.2888}.
  \item Z. Li and M.-C. Lai, \emph{The immersed interface method for the Navier--Stokes equations with singular forces}, \emph{Journal of Computational Physics} 171 (2001), 822--842. \href{https://doi.org/10.1006/jcph.2001.6813}{doi:10.1006/jcph.2001.6813}.
\end{itemize}
\begin{punchbox}
\small Free copies of the second and third papers are on PubMed Central (\href{https://pmc.ncbi.nlm.nih.gov/articles/PMC7531444/}{PMC7531444} and \href{https://pmc.ncbi.nlm.nih.gov/articles/PMC5650596/}{PMC5650596}). The repository's \texttt{papers/README.md} lists links for all of them.
\end{punchbox}
\end{frame}

\begin{frame}{References (continued)}
\footnotesize
\begin{itemize}\itemsep5pt
  \item C. S. Peskin, \emph{Flow patterns around heart valves: a numerical method}, \emph{Journal of Computational Physics} 10 (1972), 252--271. \href{https://doi.org/10.1016/0021-9991(72)90065-4}{doi:10.1016/0021-9991(72)90065-4}.
  \item C. S. Peskin, \emph{Numerical analysis of blood flow in the heart}, \emph{Journal of Computational Physics} 25 (1977), 220--252. \href{https://doi.org/10.1016/0021-9991(77)90100-0}{doi:10.1016/0021-9991(77)90100-0}.
  \item R. J. LeVeque and Z. Li, \emph{The immersed interface method for elliptic equations with discontinuous coefficients and singular sources}, \emph{SIAM Journal on Numerical Analysis} 31 (1994), 1019--1044. \href{https://doi.org/10.1137/0731054}{doi:10.1137/0731054}.
  \item B. E. Griffith, R. D. Hornung, D. M. McQueen, and C. S. Peskin, \emph{An adaptive, formally second order accurate version of the immersed boundary method}, \emph{Journal of Computational Physics} 223 (2007), 10--49. (IBAMR.)
  \item N. A. Battista, W. C. Strickland, and L. A. Miller, \emph{IB2d: a Python and MATLAB implementation of the immersed boundary method}, \emph{Bioinspiration \& Biomimetics} 12 (2017), 036003. \href{https://doi.org/10.1088/1748-3190/aa5e08}{doi:10.1088/1748-3190/aa5e08}.
  \item Wikimedia Commons, figure ``ImmersedBoundary grid,'' used in the Wikipedia article ``Immersed boundary method,'' \url{https://en.wikipedia.org/wiki/Immersed_boundary_method}. Source of the schematics \texttt{ImmersedBoundary\_grid.png} and \texttt{ImmersedBoundary\_grid2.png}; licence to be confirmed.
\end{itemize}
\end{frame}

\end{document}
